
\documentclass[oneside, 10pt, a4paper]{article}
\usepackage[usenames]{color}\color[rgb]{0.000,0.000,0.000} %used for font color
\usepackage{color}

\usepackage{algorithmicx}
\usepackage{graphicx} 
\usepackage{array}
\usepackage[table]{xcolor}
\usepackage{subfigure}
\usepackage{url}
\usepackage{comment}
\usepackage{fancyhdr}
\usepackage{float}
\usepackage{multicol}
\usepackage{amsmath}
\usepackage{cite}
\usepackage{multirow}
\usepackage{adjustbox}
\usepackage{fancybox}
\usepackage{paralist}
\usepackage{setspace}  
\usepackage{amssymb}
\usepackage{textcomp}
\usepackage{booktabs}
\newcommand{\ra}[1]{\renewcommand{\arraystretch}{#1}}
\newcounter{tablerow}
\usepackage[hidelinks]{hyperref}
\usepackage{cleveref}
\usepackage{pifont}
\newcommand{\crefrangeconjunction}{--}
\crefformat{tablerow}{#2#1#3}
\crefmultiformat{tablerow}{#2#1#3}%
{ و~#2#1#3}%
{, #2#1#3}%
{ و~#2#1#3}

\renewcommand{\labelitemi}{\scriptsize$\blacksquare$} 

%\usepackage[style=ieee]{biblatex}
%\bibliography{resources}

%\doublespacing
\usepackage{endnotes}
\usepackage{tikz}
\usetikzlibrary{mindmap,trees}
\usetikzlibrary{decorations.text}
%\usepackage[top=29mm, bottom=29mm, left=19mm, right=19mm]{geometry}
\usepackage[top=1.1in, bottom=1.1in, left=0.75in, right=0.75in]{geometry}
\usepackage{xepersian}
\NewDocumentCommand{\rot}{O{60} O{1em} m}{\makebox[#2][l]{\rotatebox{#1}{#3}}}%
\newtheorem{definition}{تعريف}
\newtheorem{theorem}{قضيه}
\defpersianfont\keywordsfont[Scale=1.2]{B Nazanin}
\defpersianfont\titlefont[Scale=1.22]{B Nazanin}
\defpersianfont\secondtitlefont[Scale=1]{B Nazanin} %1.4
\defpersianfont\secondtitlelatinfont[Scale=1.32]{Times New Roman}

\settextfont[Scale=1.1]{B Nazanin}
%\settextfont[Scale=1.1]{XB Niloofar}
\setlatintextfont[Scale=0.9]{Times New Roman}\setdigitfont[Scale=1]{B Nazanin}
%\setromantextfont[Scale=1.2]{XB Niloofar} 

\defpersianfont\sectionsfont[Scale=0.9]{B Nazanin}
\defpersianfont\captionsfont[Scale=1.0]{B Nazanin}

%\let\oldcaption=\caption
%\renewcommand{\caption}[1]{{\bf \captionsfont \oldcaption{#1}}}

\usepackage{sectsty}
\allsectionsfont{\sectionsfont}


\linespread{1.4} %line space
%\setstretch{1}

\setlength{\parindent}{5pt} %The indent of the paragraph first line
\setlength{\parskip}{5pt} %Befor paragraph space

\setlength{\columnsep}{6mm}
%\let\footnote=\endnote
%\renewcommand{\notesname}{واژه‌نامه}
\renewcommand\refname{\raggedleft \rl{منابع و مراجع}}
\newcommand{\cmark}{\ding{51}}%
\newcommand{\xmark}{\ding{55}}%
\def\InE{\begingroup\beginL\latinfont}
\def\EnE{\endL\endgroup} 
\def\InF{\begingroup\beginR\persianfont}
\def\EnF{\endR\endgroup}
\newcommand{\IE}[1]{\lr{#1}}
\newcommand{\IP}[1]{\rl{#1}} 
\newcommand{\IF}[1]{\rl{#1}}
\def\persiandash{\rl{-}}
%\defpersianfont\myfont{XB Niloofar}
\newcommand{\efootnote}[1]{\endnote{\lr{#1}}}
%\newcommand{\ecfootnote}[1]{}

\let\oldemph=\emph
\renewcommand{\emph}[1]{''\oldemph{#1}``}
\definecolor{mycolor7}{RGB}{204,255,102}
\definecolor{mycolor2}{RGB}{0,255,102}
\definecolor{mycolor3}{RGB}{255,255,153}
\definecolor{mycolor4}{RGB}{204,153,255}
\definecolor{mycolor5}{RGB}{255,204,153}
\definecolor{mycolor6}{RGB}{255,153,153}
\definecolor{mycolor}{RGB}{102,255,255}

%\setlength{\textheight}{24cm}
%\setlength{\oddsidemargin}{0cm}
%\addtolength{\topmargin}{-12mm}

%%%%%%%%%%%Circular%%%%%%%%%%
%\newcommand*{\mytextstyle}{\sffamily\Large\bfseries\color{black!85}}
\newcommand{\arcarrow}[8]{%
	% inner radius, middle radius, outer radius, start angle,
	% end angle, tip protusion angle, options, text
	\pgfmathsetmacro{\rin}{#1}
	\pgfmathsetmacro{\rmid}{#2}
	\pgfmathsetmacro{\rout}{#3}
	\pgfmathsetmacro{\astart}{#4}
	\pgfmathsetmacro{\aend}{#5}
	\pgfmathsetmacro{\atip}{#6}
	\fill[#7] (\astart:\rin) arc (\astart:\aend:\rin)
	-- (\aend+\atip:\rmid) -- (\aend:\rout) arc (\aend:\astart:\rout)
	-- (\astart+\atip:\rmid) -- cycle;
	\path[font = \sffamily, decoration = {text along path, text = {|\large\color{black}|#8},
		text align = {align = center}, raise = -0.5ex}, decorate]
	(\astart+\atip:\rmid) arc (\astart+\atip:\aend+\atip:\rmid);
}
%%%%%%%%%%%Circular%%%%%%%%%%


%\renewcommand{\thechapter}{\tartibi{chapter}}
\renewcommand{\thesection}{\arabic{section}}
\renewcommand{\thesubsection}{\arabic{section}-\arabic{subsection}}
\renewcommand{\thesubsubsection}{\arabic{section}-\arabic{subsection}-\arabic{subsubsection}}

\newcommand{\authormajor}{----}
\newcommand{\thesistitlepersian}{-----}
\newcommand{\thesistitlelatin}{Remote Attestation Technique for Industrial Control Systems}
\newcommand{\authorname}{----}
\newcommand{\studentid}{97213252}
%\footnote{
%شماره‌ي دانشجويي: 97213252، آدرس:
%\lr{mssalehi@ce.sharif.edu}
%} 

\newcommand{\advisorprofname}{-----} %persian:   استاد راهنما
\newcommand{\examinerprofname}{-----} %persian:   استاد ممتحن داخلي يا استاد مشاور

\renewcommand{\maketitle}{
\begin{titlepage}
\linespread{1.6} %line space

%\thispagestyle{fancy}

\begin{center}
\vspace*{-1cm}
\includegraphics[width=1.4in]{img/SharifLogo}\\[0.0cm]

\textbf{\titlefont
دانشگاه صنعتي شريف
\\
دانشکده مهندسي کامپيوتر
\\
سمينار کارشناسي ارشد گرايش
\authormajor
\\ [1.7cm]
}

\linespread{1.6} %line space

%\textbf{\titlefont
\settextfont[Scale=1.4]{B Nazanin}{\setlatintextfont[Scale=1.2]{Times New Roman} 
\textbf{
عنوان:
\\
\thesistitlepersian
\\
\lr{\thesistitlelatin}
\\ [2.0cm]
}
}
\linespread{1.6} %line space

\settextfont[Scale=1.4]{B Nazanin}{\setlatintextfont[Scale=1.2]{Times New Roman} 
	\textbf{
نگارش:
\\
\authorname
\\
\studentid
\\ [1.1cm]
استاد راهنما:
\\
\advisorprofname
\\ [1.0cm]
استاد ممتحن داخلي:
\\
\examinerprofname
\\ [4.0cm]
}
}

{\titlefont \textbf{
بهمن 98}
}\\[0.5cm]

\end{center}

\end{titlepage}
\setcounter{page}{1}
}


\begin{document} 
\maketitle 

\renewcommand*\thesection{\arabic{section}}
\textbf{\large{چکیده:}}
\label{sec:abstract}
لورم ایپسوم متن ساختگی با تولید سادگی نامفهوم از صنعت چاپ و با استفاده از طراحان گرافیک است. چاپگرها و متون بلکه روزنامه و مجله در ستون و سطرآنچنان که لازم است و برای شرایط فعلی تکنولوژی مورد نیاز و کاربردهای متنوع با هدف بهبود ابزارهای کاربردی می باشد. کتابهای زیادی در شصت و سه درصد گذشته، حال و آینده شناخت فراوان جامعه و متخصصان را می طلبد تا با نرم افزارها شناخت بیشتری را برای طراحان رایانه ای علی الخصوص طراحان خلاقی و فرهنگ پیشرو در زبان فارسی ایجاد کرد. در این صورت می توان امید داشت که تمام و دشواری موجود در ارائه راهکارها و شرایط سخت تایپ به پایان رسد وزمان مورد نیاز شامل حروفچینی دستاوردهای اصلی و جوابگوی سوالات پیوسته اهل دنیای موجود طراحی اساسا مورد استفاده قرار گیرد.


\noindent\keywordsfont\textbf{واژه‌های کلیدی:}
سامانه‌های بحرانی-مختلط، بسترهای چندهسته‌ای، تحمل‌پذیری اشکال، مدیریت دما

\section{مقدمه}
\label{sec:introduction}
لورم ایپسوم متن ساختگی با تولید سادگی نامفهوم از صنعت چاپ و با استفاده از طراحان گرافیک است. چاپگرها و متون بلکه روزنامه و مجله در ستون و سطرآنچنان که لازم است و برای شرایط فعلی تکنولوژی مورد نیاز و کاربردهای متنوع با هدف بهبود ابزارهای کاربردی می باشد. کتابهای زیادی در شصت و سه درصد گذشته، حال و آینده شناخت فراوان جامعه و متخصصان را می طلبد تا با نرم افزارها شناخت بیشتری را برای طراحان رایانه ای علی الخصوص طراحان خلاقی و فرهنگ پیشرو در زبان فارسی ایجاد کرد. در این صورت می توان امید داشت که تمام و دشواری موجود در ارائه راهکارها و شرایط سخت تایپ به پایان رسد وزمان مورد نیاز شامل حروفچینی دستاوردهای اصلی و جوابگوی سوالات پیوسته اهل دنیای موجود طراحی اساسا مورد استفاده قرار گیرد.

لورم ایپسوم متن ساختگی با تولید سادگی نامفهوم از صنعت چاپ و با استفاده از طراحان گرافیک است. چاپگرها و متون بلکه روزنامه و مجله در ستون و سطرآنچنان که لازم است و برای شرایط فعلی تکنولوژی مورد نیاز و کاربردهای متنوع با هدف بهبود ابزارهای کاربردی می باشد. کتابهای زیادی در شصت و سه درصد گذشته، حال و آینده شناخت فراوان جامعه و متخصصان را می طلبد تا با نرم افزارها شناخت بیشتری را برای طراحان رایانه ای علی الخصوص طراحان خلاقی و فرهنگ پیشرو در زبان فارسی ایجاد کرد. در این صورت می توان امید داشت که تمام و دشواری موجود در ارائه راهکارها و شرایط سخت تایپ به پایان رسد وزمان مورد نیاز شامل حروفچینی دستاوردهای اصلی و جوابگوی سوالات پیوسته اهل دنیای موجود طراحی اساسا مورد استفاده قرار گیرد.


%\newline\newline

 



%=========== Regular chapters ===============
%\cleardoublepage % terminates the current paragraph and page, the same way as a report document.
\setlatintextfont[Scale=0.8]{Times New Roman} 
\linespread{.8} %line space
\renewcommand{\url}[1]{\underline{#1}}

%\bibliographystyle{IEEEtran}
%\bibliographystyle{is-unsrt}
\bibliographystyle{ieeetr-fa}
%\bibliographystyle{amsplain}
\newpage
\latin
\bibliography{resources}

\persian
\linespread{1} %line space

%\theendnotes
\section*{واژه‌نامه}
{
	\renewcommand{\section}[2]{} % disable the automatic chapter
	\begin{RTL}
		\begin{multicols}{3}
			\theendnotes
		\end{multicols}
	\end{RTL}
}
\end{document}
\end{document}
